%%% appendix.tex — Приложения (полный исходный код системы NetStego)
\newpage
\section*{ПРИЛОЖЕНИЕ А. Модуль обнаружения стеганографических каналов}
\addcontentsline{toc}{section}{ПРИЛОЖЕНИЕ А. Модуль обнаружения стеганографических каналов}

Модуль \texttt{detection/analyzer.py} является центральным компонентом
подсистемы обнаружения скрытых каналов. Оркестрирует работу статистических
тестов, машинного обучения и сигнатурного анализа, вычисляет взвешенную
оценку уверенности и формирует итоговый отчёт.

\begin{lstlisting}[style=appendixcode]
"""Detection orchestrator: combines statistical, ML, and signature methods.

Analyzes PCAP files or live packet captures for signs of
steganographic covert channels.
"""

import json
from dataclasses import dataclass, field
from pathlib import Path

from rich.console import Console
from rich.table import Table
from scapy.layers.dns import DNSQR
from scapy.layers.inet import ICMP, IP, TCP
from scapy.packet import Raw
from scapy.utils import rdpcap

from netstego.detection.ml_detector import (
    MLResult, detect_isolation_forest, extract_features,
)
from netstego.detection.signatures import (
    SignatureMatch, detect_signatures,
)
from netstego.detection.statistical import (
    StatResult, chi_squared_uniformity, ipd_regularity,
    ks_test_ipd, shannon_entropy,
)

console = Console()


@dataclass
class AnalysisReport:
    """Complete analysis report."""

    source: str
    total_packets: int
    statistical_results: list[StatResult] = field(
        default_factory=list
    )
    ml_result: MLResult | None = None
    signature_matches: list[SignatureMatch] = field(
        default_factory=list
    )
    is_suspicious: bool = False
    confidence: float = 0.0

    def to_dict(self) -> dict:
        """Serialize report to dictionary."""
        proto_sigs = [
            m for m in self.signature_matches
            if m.field != "payload_content"
        ]
        content_sigs = [
            m for m in self.signature_matches
            if m.field == "payload_content"
        ]
        return {
            "source": self.source,
            "total_packets": self.total_packets,
            "is_suspicious": self.is_suspicious,
            "confidence": self.confidence,
            "statistical": [
                {
                    "method": r.method,
                    "metric": r.metric,
                    "score": r.score,
                    "p_value": r.p_value,
                    "is_anomaly": r.is_anomaly,
                    "details": r.details,
                }
                for r in self.statistical_results
            ],
            "ml": {
                "method": self.ml_result.method,
                "anomaly_ratio": self.ml_result.anomaly_ratio,
                "is_anomaly": self.ml_result.is_anomaly,
                "details": self.ml_result.details,
            }
            if self.ml_result
            else None,
            "signatures": [
                {
                    "packet_index": m.packet_index,
                    "pattern": m.pattern,
                    "confidence": m.confidence,
                    "details": m.details,
                }
                for m in proto_sigs
            ],
            "content_analysis": [
                {
                    "packet_index": m.packet_index,
                    "pattern": m.pattern,
                    "confidence": m.confidence,
                    "details": m.details,
                }
                for m in content_sigs
            ],
        }


def _extract_packet_data(packets: list) -> dict:
    """Extract field values, timestamps, and payloads from packets.

    Returns dict with keys: ip_ids, tcp_seqs, timestamps,
    payloads, ipds, sizes, tcp_tsvals, dns_qnames,
    icmp_payload_sizes.
    """
    ip_ids = []
    tcp_seqs = []
    tcp_tsvals = []
    dns_qnames = []
    icmp_payload_sizes = []
    timestamps = []
    payloads = []
    sizes = []

    for pkt in packets:
        if pkt.haslayer(IP):
            ip_ids.append(pkt[IP].id)
        if pkt.haslayer(TCP):
            tcp_seqs.append(pkt[TCP].seq)
            for opt_name, opt_val in pkt[TCP].options:
                if opt_name == "Timestamp":
                    tcp_tsvals.append(opt_val[0])
                    break
        if pkt.haslayer(DNSQR):
            qname = pkt[DNSQR].qname
            if isinstance(qname, bytes):
                qname = qname.decode(
                    "ascii", errors="ignore"
                )
            dns_qnames.append(qname.rstrip("."))
        if pkt.haslayer(ICMP) and pkt.haslayer(Raw):
            icmp_payload_sizes.append(
                len(bytes(pkt[Raw].load))
            )
        ts = (
            float(pkt.time)
            if hasattr(pkt, "time")
            else 0.0
        )
        timestamps.append(ts)
        if pkt.haslayer(Raw):
            payloads.append(bytes(pkt[Raw].load))
        else:
            payloads.append(b"")
        sizes.append(len(pkt))

    ipds = [
        timestamps[i] - timestamps[i - 1]
        for i in range(1, len(timestamps))
    ]
    return {
        "ip_ids": ip_ids,
        "tcp_seqs": tcp_seqs,
        "tcp_tsvals": tcp_tsvals,
        "dns_qnames": dns_qnames,
        "icmp_payload_sizes": icmp_payload_sizes,
        "timestamps": timestamps,
        "payloads": payloads,
        "ipds": ipds,
        "sizes": sizes,
    }


def _ip_id_sequentiality(ip_ids: list[int]) -> StatResult:
    """Check if IP ID values follow a sequential pattern.

    Normal Windows traffic: IP ID increments by 1-3.
    Stego: arbitrary values that break sequential patterns.
    """
    if len(ip_ids) < 3:
        return StatResult(
            method="ip_id_sequentiality",
            metric="sequential_ratio",
            score=0.0, p_value=None,
            threshold=0.0, is_anomaly=False,
            details="Insufficient data",
        )
    sequential_count = 0
    for i in range(1, len(ip_ids)):
        delta = (ip_ids[i] - ip_ids[i - 1]) % 65536
        if 1 <= delta <= 10:
            sequential_count += 1
    ratio = sequential_count / (len(ip_ids) - 1)
    threshold = 0.3
    return StatResult(
        method="ip_id_sequentiality",
        metric="sequential_ratio",
        score=ratio, p_value=None,
        threshold=threshold,
        is_anomaly=ratio < threshold,
        details=(
            f"sequential_pairs="
            f"{sequential_count}/{len(ip_ids)-1}, "
            f"ratio={ratio:.4f}"
        ),
    )


def _tsval_monotonicity(
    tsvals: list[int],
) -> StatResult:
    """Check if TCP TSval is monotonically increasing.

    Normal: monotonic with consistent step.
    Stego: arbitrary data breaking monotonicity.
    """
    violations = 0
    for i in range(1, len(tsvals)):
        if tsvals[i] <= tsvals[i - 1]:
            violations += 1
    ratio = (
        violations / (len(tsvals) - 1)
        if len(tsvals) > 1
        else 0.0
    )
    threshold = 0.2
    return StatResult(
        method="tsval_monotonicity",
        metric="violation_ratio",
        score=ratio, p_value=None,
        threshold=threshold,
        is_anomaly=ratio > threshold,
        details=(
            f"violations={violations}/"
            f"{len(tsvals)-1}, ratio={ratio:.4f}"
        ),
    )


def _dns_qname_entropy(
    qnames: list[str],
) -> StatResult:
    """Check entropy of DNS QNAME characters.

    Normal DNS: readable names, low entropy (~3.5-4.0).
    Base32 stego: random, high entropy (~4.5-5.0).
    """
    import math

    all_chars = "".join(qnames)
    if not all_chars:
        return StatResult(
            method="dns_qname_entropy",
            metric="char_entropy",
            score=0.0, p_value=None,
            threshold=0.0, is_anomaly=False,
            details="No data",
        )
    freq: dict[str, int] = {}
    for c in all_chars:
        freq[c] = freq.get(c, 0) + 1
    n = len(all_chars)
    entropy = -sum(
        (cnt / n) * math.log2(cnt / n)
        for cnt in freq.values()
    )
    avg_len = sum(len(q) for q in qnames) / len(qnames)
    threshold_entropy = 4.3
    threshold_len = 25
    is_anomaly = (
        entropy > threshold_entropy
        and avg_len > threshold_len
    )
    return StatResult(
        method="dns_qname_entropy",
        metric="char_entropy",
        score=entropy, p_value=None,
        threshold=threshold_entropy,
        is_anomaly=is_anomaly,
        details=(
            f"H={entropy:.4f}, "
            f"avg_qname_len={avg_len:.1f}, "
            f"unique_chars={len(freq)}"
        ),
    )


def _icmp_payload_anomaly(
    payload_sizes: list[int],
) -> StatResult:
    """Check if ICMP payload sizes are abnormal.

    Normal ping: 32 bytes (Win) or 56 bytes (Linux).
    Stego: often much larger (up to 1472 bytes).
    """
    avg_size = sum(payload_sizes) / len(payload_sizes)
    max_size = max(payload_sizes)
    threshold = 100
    is_anomaly = avg_size > threshold
    return StatResult(
        method="icmp_payload_size",
        metric="avg_payload_bytes",
        score=avg_size, p_value=None,
        threshold=float(threshold),
        is_anomaly=is_anomaly,
        details=(
            f"avg={avg_size:.0f}, max={max_size}, "
            f"count={len(payload_sizes)}"
        ),
    )


def _payload_byte_entropy(
    payloads: list[bytes],
) -> StatResult:
    """Check byte-level entropy of packet payloads.

    Encrypted stego data: entropy close to 8.0 bits/byte.
    Normal ICMP: repeating pattern, much lower entropy.
    """
    import math

    all_bytes = b"".join(payloads)
    if not all_bytes:
        return StatResult(
            method="payload_byte_entropy",
            metric="bits_per_byte",
            score=0.0, p_value=None,
            threshold=0.0, is_anomaly=False,
            details="No data",
        )
    freq = [0] * 256
    for b in all_bytes:
        freq[b] += 1
    n = len(all_bytes)
    entropy = -sum(
        (f / n) * math.log2(f / n) for f in freq if f > 0
    )
    threshold = 7.0
    return StatResult(
        method="payload_byte_entropy",
        metric="bits_per_byte",
        score=entropy, p_value=None,
        threshold=threshold,
        is_anomaly=entropy > threshold,
        details=(
            f"H={entropy:.4f} bits/byte, "
            f"total_bytes={n}"
        ),
    )


def _detect_dns_stego_signatures(
    qnames: list[str],
) -> list:
    """Detect stego patterns in DNS QNAME labels."""
    import re
    from netstego.detection.signatures import (
        SignatureMatch,
    )

    matches = []
    base32_pattern = re.compile(r"^[a-z2-7]{20,}$")
    for i, qname in enumerate(qnames):
        labels = qname.split(".")
        for label in labels[:-2]:
            if (
                base32_pattern.match(label)
                and len(label) > 20
            ):
                matches.append(SignatureMatch(
                    packet_index=i,
                    field="dns_qname",
                    pattern="BASE32_SUBDOMAIN",
                    confidence=0.85,
                    details=(
                        f"Base32-like label "
                        f"({len(label)} chars): "
                        f"{label[:30]}..."
                    ),
                ))
                break
        if any(l == "term" for l in labels):
            matches.append(SignatureMatch(
                packet_index=i,
                field="dns_qname",
                pattern="DNS_TERM_MARKER",
                confidence=0.9,
                details=(
                    f"Termination marker "
                    f"in QNAME: {qname}"
                ),
            ))
    return matches


def analyze_packets(
    packets: list,
    source: str = "unknown",
    methods: list[str] | None = None,
    alpha: float = 0.05,
    ml_contamination: float = 0.05,
) -> AnalysisReport:
    """Run detection analysis on a list of packets.

    Args:
        packets: List of Scapy packets.
        source: Source description (file path or iface).
        methods: Methods to run: 'statistical',
                 'ml', 'signatures'. Default: all.
        alpha: Significance level for stat tests.
        ml_contamination: IsolationForest parameter.

    Returns:
        AnalysisReport with all results.
    """
    if methods is None:
        methods = ["statistical", "ml", "signatures"]

    report = AnalysisReport(
        source=source, total_packets=len(packets)
    )
    data = _extract_packet_data(packets)

    # Statistical tests
    if "statistical" in methods:
        if (
            data["ip_ids"]
            and len(data["ip_ids"]) >= 10
        ):
            report.statistical_results.append(
                _ip_id_sequentiality(data["ip_ids"])
            )
        if (
            data["tcp_seqs"]
            and len(data["tcp_seqs"]) >= 10
        ):
            report.statistical_results.append(
                chi_squared_uniformity(
                    data["tcp_seqs"], alpha
                )
            )
        if data["ipds"]:
            report.statistical_results.append(
                ipd_regularity(data["ipds"])
            )
        if (
            data["tcp_tsvals"]
            and len(data["tcp_tsvals"]) >= 3
        ):
            report.statistical_results.append(
                _tsval_monotonicity(data["tcp_tsvals"])
            )
        if (
            data["dns_qnames"]
            and len(data["dns_qnames"]) >= 3
        ):
            report.statistical_results.append(
                _dns_qname_entropy(data["dns_qnames"])
            )
        if (
            data["icmp_payload_sizes"]
            and len(data["icmp_payload_sizes"]) >= 1
        ):
            report.statistical_results.append(
                _icmp_payload_anomaly(
                    data["icmp_payload_sizes"]
                )
            )
        icmp_payloads = [
            p for p in data["payloads"] if len(p) > 40
        ]
        if icmp_payloads:
            report.statistical_results.append(
                _payload_byte_entropy(icmp_payloads)
            )

    # ML detection
    if "ml" in methods:
        field_values = (
            data["ip_ids"]
            or data["tcp_seqs"]
            or [0] * len(packets)
        )
        features = extract_features(
            field_values, data["timestamps"],
            data["sizes"], data["payloads"],
        )
        if len(features) >= 10:
            report.ml_result = detect_isolation_forest(
                features, ml_contamination
            )

    # Signature detection
    if "signatures" in methods:
        report.signature_matches = detect_signatures(
            data["payloads"]
        )
        if data["dns_qnames"]:
            report.signature_matches.extend(
                _detect_dns_stego_signatures(
                    data["dns_qnames"]
                )
            )

    # Compute overall suspicion
    import math as _math

    valid_stats = [
        r
        for r in report.statistical_results
        if "Insufficient" not in r.details
        and "No data" not in r.details
        and (
            r.p_value is None
            or not _math.isnan(r.p_value)
        )
        and "bins=1" not in r.details
    ]
    anomaly_count = sum(
        1 for r in valid_stats if r.is_anomaly
    )
    total_tests = len(valid_stats)
    stat_ratio = (
        anomaly_count / total_tests
        if total_tests > 0
        else 0.0
    )

    strong_indicators = {
        r.method for r in valid_stats if r.is_anomaly
    } & {
        "tsval_monotonicity",
        "dns_qname_entropy",
        "icmp_payload_size",
        "payload_byte_entropy",
        "ip_id_sequentiality",
    }

    ml_suspicious = (
        report.ml_result.is_anomaly
        if report.ml_result
        else False
    )
    sig_count = len(report.signature_matches)

    confidence = 0.0
    if stat_ratio >= 0.5 or len(strong_indicators) >= 1:
        confidence += 0.4
    if ml_suspicious:
        confidence += 0.3
    if sig_count > 0:
        confidence += 0.3

    report.is_suspicious = confidence >= 0.4
    report.confidence = confidence

    return report


def analyze_pcap(
    pcap_path: Path,
    methods: list[str] | None = None,
    alpha: float = 0.05,
    ml_contamination: float = 0.05,
) -> AnalysisReport:
    """Analyze a PCAP file for covert channels."""
    packets = rdpcap(str(pcap_path))
    return analyze_packets(
        list(packets),
        source=str(pcap_path),
        methods=methods,
        alpha=alpha,
        ml_contamination=ml_contamination,
    )


def print_report(report: AnalysisReport) -> None:
    """Print analysis report to console using rich."""
    verdict = (
        "[bold red]SUSPICIOUS[/bold red]"
        if report.is_suspicious
        else "[bold green]CLEAN[/bold green]"
    )
    console.print(f"\n{'='*60}")
    console.print(
        f"[bold]Detection Report:[/bold] "
        f"{report.source}"
    )
    console.print(
        f"Packets analyzed: {report.total_packets}"
    )
    console.print(
        f"Verdict: {verdict} "
        f"(confidence: {report.confidence:.0%})"
    )
    console.print(f"{'='*60}")

    if report.statistical_results:
        table = Table(title="Statistical Tests")
        table.add_column("Method")
        table.add_column("Metric")
        table.add_column("Score")
        table.add_column("P-value")
        table.add_column("Anomaly")
        for r in report.statistical_results:
            style = "red" if r.is_anomaly else "green"
            table.add_row(
                r.method, r.metric,
                f"{r.score:.4f}",
                (
                    f"{r.p_value:.6f}"
                    if r.p_value is not None
                    else "N/A"
                ),
                f"[{style}]"
                f"{'YES' if r.is_anomaly else 'NO'}"
                f"[/{style}]",
            )
        console.print(table)

    if report.ml_result:
        ml = report.ml_result
        style = "red" if ml.is_anomaly else "green"
        console.print(
            f"\n[bold]ML Detection "
            f"({ml.method}):[/bold]"
        )
        console.print(
            f"  Anomaly ratio: "
            f"[{style}]{ml.anomaly_ratio:.2%}"
            f"[/{style}]"
        )
        console.print(f"  Details: {ml.details}")

    proto_sigs = [
        m for m in report.signature_matches
        if m.field != "payload_content"
    ]
    content_sigs = [
        m for m in report.signature_matches
        if m.field == "payload_content"
    ]
    if proto_sigs:
        console.print(
            f"\n[bold red]Protocol signatures "
            f"found: {len(proto_sigs)}[/bold red]"
        )
        for m in proto_sigs[:10]:
            console.print(
                f"  Packet #{m.packet_index}: "
                f"{m.pattern} -- {m.details}"
            )
    if content_sigs:
        console.print(
            f"\n[bold red]Malicious content "
            f"indicators: {len(content_sigs)}"
            f"[/bold red]"
        )
        for m in content_sigs[:10]:
            console.print(
                f"  Packet #{m.packet_index}: "
                f"[bold]{m.pattern}[/bold] "
                f"(confidence: {m.confidence:.0%})"
                f" -- {m.details}"
            )
\end{lstlisting}

\newpage
\section*{ПРИЛОЖЕНИЕ Б. Модули статистического анализа, машинного обучения и сигнатурного обнаружения}
\addcontentsline{toc}{section}{ПРИЛОЖЕНИЕ Б. Модули статистического анализа, машинного обучения и сигнатурного обнаружения}

\subsection*{Б.1 Статистические методы обнаружения (\texttt{detection/statistical.py})}

Модуль реализует четыре статистических теста: критерий хи-квадрат
для проверки равномерности распределения значений полей,
критерий Колмогорова--Смирнова для анализа межпакетных задержек,
вычисление энтропии Шеннона и анализ регулярности задержек.

\begin{lstlisting}[style=appendixcode]
"""Statistical detection methods for covert channel
analysis.

Implements:
    - Chi-squared test for field value distributions
    - Kolmogorov-Smirnov test for inter-packet delay
    - Shannon entropy calculation
    - Inter-packet delay analysis
"""

import math
from dataclasses import dataclass

import numpy as np
from scipy import stats as sp_stats


@dataclass
class StatResult:
    """Result of a statistical test."""

    method: str
    metric: str
    score: float
    p_value: float | None
    threshold: float
    is_anomaly: bool
    details: str = ""


def chi_squared_uniformity(
    values: list[int], alpha: float = 0.05,
) -> StatResult:
    """Chi-squared test for uniformity of field values.

    Tests whether the distribution of values deviates
    significantly from a uniform distribution, which
    may indicate steganographic encoding.
    """
    if len(values) < 10:
        return StatResult(
            method="chi_squared",
            metric="uniformity",
            score=0.0, p_value=1.0,
            threshold=alpha, is_anomaly=False,
            details="Insufficient data (< 10 samples)",
        )

    arr = np.array(values)
    n_bins = min(256, len(set(values)))
    observed, bin_edges = np.histogram(
        arr, bins=n_bins
    )
    expected = np.full_like(
        observed,
        fill_value=len(values) / n_bins,
        dtype=float,
    )
    chi2, p_value = sp_stats.chisquare(
        observed, f_exp=expected
    )
    return StatResult(
        method="chi_squared",
        metric="uniformity",
        score=float(chi2),
        p_value=float(p_value),
        threshold=alpha,
        is_anomaly=p_value < alpha,
        details=(
            f"chi2={chi2:.4f}, p={p_value:.6f}, "
            f"bins={n_bins}"
        ),
    )


def ks_test_ipd(
    delays: list[float], alpha: float = 0.05,
) -> StatResult:
    """Kolmogorov-Smirnov test for IPD distribution.

    Tests whether IPD follows an exponential
    distribution (typical for normal traffic).
    Steganographic traffic has more regular timing.
    """
    if len(delays) < 10:
        return StatResult(
            method="ks_test",
            metric="ipd_distribution",
            score=0.0, p_value=1.0,
            threshold=alpha, is_anomaly=False,
            details="Insufficient data (< 10 samples)",
        )

    arr = np.array(delays)
    arr = arr[arr > 0]
    if len(arr) < 5:
        return StatResult(
            method="ks_test",
            metric="ipd_distribution",
            score=0.0, p_value=1.0,
            threshold=alpha, is_anomaly=False,
            details="Insufficient positive delays",
        )

    mean_delay = float(np.mean(arr))
    stat, p_value = sp_stats.kstest(
        arr, "expon", args=(0, mean_delay)
    )
    return StatResult(
        method="ks_test",
        metric="ipd_distribution",
        score=float(stat),
        p_value=float(p_value),
        threshold=alpha,
        is_anomaly=p_value < alpha,
        details=(
            f"D={stat:.4f}, p={p_value:.6f}, "
            f"mean_ipd={mean_delay:.4f}s"
        ),
    )


def shannon_entropy(
    values: list[int],
) -> StatResult:
    """Calculate Shannon entropy of field values.

    High entropy close to theoretical maximum may
    indicate steganographic encoding.
    """
    if not values:
        return StatResult(
            method="entropy",
            metric="shannon",
            score=0.0, p_value=None,
            threshold=0.0, is_anomaly=False,
            details="No data",
        )

    n = len(values)
    freq: dict[int, int] = {}
    for v in values:
        freq[v] = freq.get(v, 0) + 1

    entropy = 0.0
    for count in freq.values():
        p = count / n
        if p > 0:
            entropy -= p * math.log2(p)

    max_entropy = math.log2(n) if n > 1 else 1.0
    normalized = (
        entropy / max_entropy
        if max_entropy > 0
        else 0.0
    )
    threshold = 0.95
    return StatResult(
        method="entropy",
        metric="shannon",
        score=normalized,
        p_value=None,
        threshold=threshold,
        is_anomaly=normalized > threshold,
        details=(
            f"H={entropy:.4f} bits, "
            f"normalized={normalized:.4f}, "
            f"unique={len(freq)}"
        ),
    )


def ipd_regularity(
    delays: list[float],
) -> StatResult:
    """Measure regularity of inter-packet delays.

    Low coefficient of variation (CV) suggests
    artificially timed traffic.
    """
    if len(delays) < 5:
        return StatResult(
            method="ipd_regularity",
            metric="coefficient_of_variation",
            score=0.0, p_value=None,
            threshold=0.0, is_anomaly=False,
            details="Insufficient data",
        )

    arr = np.array(delays)
    mean = float(np.mean(arr))
    std = float(np.std(arr))
    cv = std / mean if mean > 0 else 0.0

    threshold = 0.1
    return StatResult(
        method="ipd_regularity",
        metric="coefficient_of_variation",
        score=cv, p_value=None,
        threshold=threshold,
        is_anomaly=cv < threshold,
        details=(
            f"CV={cv:.4f}, mean={mean:.4f}s, "
            f"std={std:.4f}s"
        ),
    )
\end{lstlisting}

\newpage
\subsection*{Б.2 Обнаружение на основе машинного обучения}
\noindent Модуль \texttt{detection/ml\_detector.py}

Модуль реализует два метода обнаружения на основе машинного обучения:
Isolation Forest (обучение без учителя) и Random Forest (обучение
с учителем). Извлекает 9-мерный вектор признаков из пакетных данных.

\begin{lstlisting}[style=appendixcode]
"""Machine learning based covert channel detection.

Uses IsolationForest for anomaly detection on packet
features. Can optionally train a RandomForest classifier
if labeled data is available.
"""

from dataclasses import dataclass

import numpy as np
from sklearn.ensemble import (
    IsolationForest, RandomForestClassifier,
)
from sklearn.preprocessing import StandardScaler


@dataclass
class MLResult:
    """Result of ML-based detection."""

    method: str
    anomaly_ratio: float
    predictions: list[int]
    confidence: float
    is_anomaly: bool
    details: str = ""


def _byte_entropy(data: bytes) -> float:
    """Shannon entropy of byte values (0-8 bits/byte)."""
    if not data:
        return 0.0
    freq = [0] * 256
    for b in data:
        freq[b] += 1
    n = len(data)
    import math
    return -sum(
        (f / n) * math.log2(f / n)
        for f in freq
        if f > 0
    )


def extract_features(
    field_values: list[int],
    timestamps: list[float],
    packet_sizes: list[int] | None = None,
    payloads: list[bytes] | None = None,
) -> np.ndarray:
    """Extract feature vectors from packet data.

    Features per packet (sliding window of 5):
        - Field value
        - Delta from previous value
        - Inter-packet delay
        - Packet size
        - Rolling mean of field values
        - Rolling std of field values
        - Rolling mean of IPD
        - Payload byte entropy (0-8 bits/byte)
        - Payload length

    Returns:
        2D numpy array of shape (n_samples, 9).
    """
    n = len(field_values)
    if n < 2:
        return np.empty((0, 9))

    values = np.array(field_values, dtype=float)
    times = np.array(timestamps, dtype=float)
    sizes = (
        np.array(packet_sizes, dtype=float)
        if packet_sizes
        else np.zeros(n)
    )

    if payloads and len(payloads) == n:
        payload_entropies = np.array(
            [_byte_entropy(p) for p in payloads]
        )
        payload_lengths = np.array(
            [len(p) for p in payloads], dtype=float
        )
    else:
        payload_entropies = np.zeros(n)
        payload_lengths = np.zeros(n)

    value_deltas = np.diff(values, prepend=values[0])
    ipd = np.diff(times, prepend=times[0])

    window = min(5, n)
    features = []
    for i in range(n):
        start = max(0, i - window + 1)
        win_vals = values[start : i + 1]
        win_ipd = ipd[start : i + 1]
        features.append([
            values[i],
            value_deltas[i],
            ipd[i],
            sizes[i],
            float(np.mean(win_vals)),
            float(np.std(win_vals)),
            float(np.mean(win_ipd)),
            payload_entropies[i],
            payload_lengths[i],
        ])

    return np.array(features)


def detect_isolation_forest(
    features: np.ndarray,
    contamination: float = 0.05,
) -> MLResult:
    """Run IsolationForest anomaly detection."""
    if len(features) < 10:
        return MLResult(
            method="isolation_forest",
            anomaly_ratio=0.0, predictions=[],
            confidence=0.0, is_anomaly=False,
            details="Insufficient data (< 10 samples)",
        )

    scaler = StandardScaler()
    scaled = scaler.fit_transform(features)

    clf = IsolationForest(
        contamination=contamination,
        random_state=42,
        n_estimators=100,
    )
    predictions = clf.fit_predict(scaled)
    scores = clf.decision_function(scaled)

    anomaly_count = int(np.sum(predictions == -1))
    anomaly_ratio = anomaly_count / len(predictions)
    mean_score = float(np.mean(scores))

    return MLResult(
        method="isolation_forest",
        anomaly_ratio=anomaly_ratio,
        predictions=predictions.tolist(),
        confidence=abs(mean_score),
        is_anomaly=anomaly_ratio > contamination * 2,
        details=(
            f"anomalies={anomaly_count}/"
            f"{len(predictions)}, "
            f"ratio={anomaly_ratio:.4f}, "
            f"mean_score={mean_score:.4f}"
        ),
    )


def train_classifier(
    features: np.ndarray,
    labels: np.ndarray,
) -> tuple[RandomForestClassifier, StandardScaler]:
    """Train a RandomForest classifier on labeled data.

    Returns:
        Tuple of (trained classifier, fitted scaler).
    """
    scaler = StandardScaler()
    scaled = scaler.fit_transform(features)

    clf = RandomForestClassifier(
        n_estimators=100,
        random_state=42,
        max_depth=10,
    )
    clf.fit(scaled, labels)
    return clf, scaler


def detect_random_forest(
    train_features: np.ndarray,
    train_labels: np.ndarray,
    test_features: np.ndarray,
    test_labels: np.ndarray | None = None,
) -> MLResult:
    """Run supervised RandomForest detection.

    Trains on labeled normal/stego data and predicts
    test samples.
    """
    if (
        len(train_features) < 10
        or len(test_features) < 2
    ):
        return MLResult(
            method="random_forest",
            anomaly_ratio=0.0, predictions=[],
            confidence=0.0, is_anomaly=False,
            details=(
                "Insufficient data for supervised "
                "classification"
            ),
        )

    clf, scaler = train_classifier(
        train_features, train_labels
    )
    scaled_test = scaler.transform(test_features)
    predictions = clf.predict(scaled_test)
    probabilities = clf.predict_proba(scaled_test)

    stego_count = int(np.sum(predictions == 1))
    stego_ratio = (
        stego_count / len(predictions)
        if len(predictions) > 0
        else 0.0
    )
    avg_confidence = float(
        np.mean(np.max(probabilities, axis=1))
    )

    details_parts = [
        f"stego={stego_count}/{len(predictions)}",
        f"ratio={stego_ratio:.4f}",
        f"avg_confidence={avg_confidence:.4f}",
    ]

    if (
        test_labels is not None
        and len(test_labels) == len(predictions)
    ):
        accuracy = float(
            np.mean(predictions == test_labels)
        )
        details_parts.append(
            f"accuracy={accuracy:.4f}"
        )

    return MLResult(
        method="random_forest",
        anomaly_ratio=stego_ratio,
        predictions=predictions.tolist(),
        confidence=avg_confidence,
        is_anomaly=stego_ratio > 0.3,
        details=", ".join(details_parts),
    )
\end{lstlisting}

\newpage
\subsection*{Б.3 Сигнатурный анализ и обнаружение вредоносных вложений}
\noindent Модуль \texttt{detection/signatures.py}

Модуль реализует сигнатурный анализ пакетов: обнаружение маркера
протокола NST, анализ паттернов длины и проверку содержимого
извлечённых данных на наличие вредоносных вложений (20~сигнатур
в~7~категориях).

\begin{lstlisting}[style=appendixcode]
"""Signature-based detection for NetStego's internal
protocol.

Scans packet payloads and field patterns for the NST
magic header and other markers of the steganographic
protocol.
"""

from dataclasses import dataclass

from netstego.core.fragmentation import MAGIC


@dataclass
class SignatureMatch:
    """A detected signature match."""

    packet_index: int
    field: str
    pattern: str
    confidence: float
    details: str = ""


def scan_for_magic(
    payloads: list[bytes],
) -> list[SignatureMatch]:
    """Scan packet payloads for NST magic bytes."""
    matches = []
    for i, payload in enumerate(payloads):
        idx = payload.find(MAGIC)
        if idx != -1:
            matches.append(SignatureMatch(
                packet_index=i,
                field="payload",
                pattern="NST_MAGIC",
                confidence=1.0,
                details=(
                    f"Magic 'NST\\x00' at offset {idx}"
                ),
            ))
    return matches


def scan_for_length_prefix_pattern(
    payloads: list[bytes],
) -> list[SignatureMatch]:
    """Detect 4-byte length prefix framing pattern.

    Checks if payloads start with a plausible 4-byte
    big-endian length followed by data of matching size.
    """
    import struct

    matches = []
    for i, payload in enumerate(payloads):
        if len(payload) < 4:
            continue
        claimed_len = struct.unpack(
            ">I", payload[:4]
        )[0]
        if 17 <= claimed_len <= len(payload) - 4:
            data_after = payload[4 : 4 + claimed_len]
            if data_after[:4] == MAGIC:
                matches.append(SignatureMatch(
                    packet_index=i,
                    field="payload",
                    pattern="LENGTH_PREFIX_MAGIC",
                    confidence=0.95,
                    details=(
                        f"Length prefix {claimed_len} "
                        f"followed by NST magic"
                    ),
                ))
    return matches


# Known malicious content signatures
_CONTENT_SIGNATURES: list[
    tuple[bytes, str, str, float]
] = [
    # (pattern, name, description, confidence)
    (b"MZ", "PE_HEADER",
     "Windows PE executable (MZ header)", 0.9),
    (b"\x7fELF", "ELF_HEADER",
     "Linux ELF executable", 0.95),
    (b"#!/bin/", "SHEBANG_SCRIPT",
     "Unix shell script (shebang)", 0.8),
    (b"#!/usr/bin/", "SHEBANG_SCRIPT",
     "Unix shell script (shebang)", 0.8),
    (b"<script", "HTML_SCRIPT",
     "Embedded HTML/JavaScript", 0.75),
    (b"powershell", "POWERSHELL_CMD",
     "PowerShell command", 0.8),
    (b"cmd /c", "CMD_EXEC",
     "Windows cmd execution", 0.8),
    (b"cmd.exe", "CMD_EXEC",
     "Windows cmd.exe reference", 0.75),
    (b"\x90\x90\x90\x90\x90\x90\x90\x90",
     "NOP_SLED",
     "NOP sled (possible shellcode)", 0.85),
    (b"\xcc\xcc\xcc\xcc", "INT3_SLED",
     "INT3 breakpoint sled", 0.8),
    (b"PK\x03\x04", "ZIP_ARCHIVE",
     "ZIP archive (may contain executables)", 0.6),
    (b"\x1f\x8b", "GZIP_ARCHIVE",
     "GZIP compressed data", 0.5),
    (b"Rar!\x1a\x07", "RAR_ARCHIVE",
     "RAR archive", 0.6),
    (b"\xd0\xcf\x11\xe0", "OLE_DOCUMENT",
     "OLE2 document (Office macro capable)", 0.7),
    (b"%PDF", "PDF_DOCUMENT",
     "PDF document (may contain JS)", 0.5),
    (b"import os", "PYTHON_IMPORT",
     "Python code with OS access", 0.65),
    (b"import subprocess", "PYTHON_SUBPROCESS",
     "Python code with subprocess", 0.75),
    (b"eval(", "EVAL_CALL",
     "Dynamic code evaluation", 0.7),
    (b"exec(", "EXEC_CALL",
     "Dynamic code execution", 0.7),
    (b"base64.b64decode", "BASE64_DECODE",
     "Base64 decoding (obfuscation)", 0.6),
]


def scan_payload_content(
    payloads: list[bytes],
) -> list[SignatureMatch]:
    """Analyze extracted payloads for signs of malicious
    content.

    Scans reassembled data for executable headers,
    scripts, shellcode patterns, and other indicators
    of malicious attachments hidden via steganographic
    channels.
    """
    matches = []
    for i, payload in enumerate(payloads):
        if len(payload) < 2:
            continue
        for pattern, name, desc, conf in (
            _CONTENT_SIGNATURES
        ):
            idx = payload.find(pattern)
            if idx != -1:
                matches.append(SignatureMatch(
                    packet_index=i,
                    field="payload_content",
                    pattern=name,
                    confidence=conf,
                    details=f"{desc} at offset {idx}",
                ))
    return matches


def detect_signatures(
    payloads: list[bytes],
) -> list[SignatureMatch]:
    """Run all signature detection methods."""
    matches = []
    matches.extend(scan_for_magic(payloads))
    matches.extend(
        scan_for_length_prefix_pattern(payloads)
    )
    matches.extend(scan_payload_content(payloads))
    return matches
\end{lstlisting}

\newpage
\section*{ПРИЛОЖЕНИЕ В. Криптографическое ядро и протокол фрагментации}
\addcontentsline{toc}{section}{ПРИЛОЖЕНИЕ В. Криптографическое ядро и протокол фрагментации}

\subsection*{В.1 Модуль шифрования (\texttt{core/crypto.py})}

Модуль реализует криптографические операции: генерацию и загрузку ключей,
деривацию ключа из пароля (Argon2id), шифрование и расшифрование
данных алгоритмом AES-256-GCM.

\begin{lstlisting}[style=appendixcode]
"""Cryptographic operations: AES-256-GCM encryption
with Argon2id KDF.

References:
    - AES-GCM: NIST SP 800-38D
    - Argon2id: RFC 9106
"""

import secrets
from pathlib import Path

from argon2.low_level import Type, hash_secret_raw
from cryptography.hazmat.primitives.ciphers.aead import (
    AESGCM,
)

KEY_SIZE = 32
NONCE_SIZE = 12
SALT_SIZE = 16


def generate_key(path: Path) -> bytes:
    """Generate a random 256-bit key and save to file."""
    key = secrets.token_bytes(KEY_SIZE)
    path.parent.mkdir(parents=True, exist_ok=True)
    path.write_bytes(key)
    return key


def load_key(path: Path) -> bytes:
    """Load a key from file.

    Raises:
        ValueError: If key file has wrong size.
    """
    data = path.read_bytes()
    if len(data) != KEY_SIZE:
        msg = (
            f"Key file must be {KEY_SIZE} bytes, "
            f"got {len(data)}"
        )
        raise ValueError(msg)
    return data


def derive_key(
    password: bytes,
    salt: bytes | None = None,
    *,
    time_cost: int = 3,
    memory_cost: int = 65536,
) -> tuple[bytes, bytes]:
    """Derive AES-256 key from password using Argon2id.

    Args:
        password: Password or passphrase bytes.
        salt: Optional salt; generated if not provided.
        time_cost: Argon2 time cost parameter.
        memory_cost: Argon2 memory cost in KiB.

    Returns:
        Tuple of (derived_key, salt).
    """
    if salt is None:
        salt = secrets.token_bytes(SALT_SIZE)
    key = hash_secret_raw(
        secret=password,
        salt=salt,
        time_cost=time_cost,
        memory_cost=memory_cost,
        parallelism=1,
        hash_len=KEY_SIZE,
        type=Type.ID,
    )
    return key, salt


def encrypt(plaintext: bytes, key: bytes) -> bytes:
    """Encrypt data with AES-256-GCM.

    Output: nonce (12 B) || ciphertext || tag (16 B)
    """
    nonce = secrets.token_bytes(NONCE_SIZE)
    aesgcm = AESGCM(key)
    ciphertext = aesgcm.encrypt(nonce, plaintext, None)
    return nonce + ciphertext


def decrypt(data: bytes, key: bytes) -> bytes:
    """Decrypt AES-256-GCM encrypted data.

    Expects: nonce (12 B) || ciphertext || tag (16 B)

    Raises:
        ValueError: If data is too short.
        InvalidTag: If authentication fails.
    """
    if len(data) < NONCE_SIZE + 16:
        msg = "Encrypted data too short"
        raise ValueError(msg)
    nonce = data[:NONCE_SIZE]
    ciphertext = data[NONCE_SIZE:]
    aesgcm = AESGCM(key)
    return aesgcm.decrypt(nonce, ciphertext, None)
\end{lstlisting}

\newpage
\subsection*{В.2 Модуль фрагментации (\texttt{core/fragmentation.py})}

Модуль реализует протокол фрагментации NST с 17-байтным заголовком чанка,
включающим магическое число, порядковый номер, общее число чанков,
контрольную сумму CRC32 и флаги.

\begin{lstlisting}[style=appendixcode]
"""Payload fragmentation with internal chunk header.

Chunk header format (17 bytes):
    Offset  Size    Field
    0       4       MAGIC: b"NST\x00"
    4       4       SEQ NUMBER (uint32 BE)
    8       4       TOTAL CHUNKS (uint32 BE)
    12      4       CRC32 (uint32 BE) of DATA only
    16      1       FLAGS (uint8)
    17      N       DATA

Flags:
    0x01 -- first chunk (DATA starts with SHA-256)
    0x02 -- last chunk
    0x04 -- NACK (retransmit request)
"""

import hashlib
import struct
import zlib

MAGIC = b"NST\x00"
HEADER_SIZE = 17

FLAG_FIRST = 0x01
FLAG_LAST = 0x02
FLAG_NACK = 0x04


def _make_header(
    seq: int, total: int, crc: int, flags: int,
) -> bytes:
    """Build a 17-byte chunk header."""
    return struct.pack(
        ">4sIIIB", MAGIC, seq, total, crc, flags
    )


def parse_header(
    chunk: bytes,
) -> tuple[int, int, int, int, bytes]:
    """Parse a chunk into (seq, total, crc32, flags, data).

    Raises:
        ValueError: If magic bytes wrong or chunk short.
    """
    if len(chunk) < HEADER_SIZE:
        msg = (
            f"Chunk too short: "
            f"{len(chunk)} < {HEADER_SIZE}"
        )
        raise ValueError(msg)
    magic, seq, total, crc, flags = struct.unpack(
        ">4sIIIB", chunk[:HEADER_SIZE]
    )
    if magic != MAGIC:
        msg = f"Invalid magic: {magic!r}"
        raise ValueError(msg)
    data = chunk[HEADER_SIZE:]
    return seq, total, crc, flags, data


def fragment(
    payload: bytes, chunk_size: int,
) -> list[bytes]:
    """Split payload into chunks with headers.

    The first chunk includes the SHA-256 hash of the
    original payload prepended to the data (32 bytes).

    Args:
        payload: Data to fragment.
        chunk_size: Max data bytes per chunk (excl. hdr).

    Raises:
        ValueError: If chunk_size < 33.
    """
    if chunk_size < 33:
        msg = (
            "chunk_size must be >= 33 to fit "
            "SHA-256 hash in first chunk"
        )
        raise ValueError(msg)

    sha256 = hashlib.sha256(payload).digest()

    first_data_limit = chunk_size - 32
    pieces: list[bytes] = []

    if len(payload) <= first_data_limit:
        pieces.append(sha256 + payload)
    else:
        pieces.append(
            sha256 + payload[:first_data_limit]
        )
        remaining = payload[first_data_limit:]
        while remaining:
            pieces.append(remaining[:chunk_size])
            remaining = remaining[chunk_size:]

    total = len(pieces)
    chunks: list[bytes] = []

    for i, data in enumerate(pieces):
        flags = 0
        if i == 0:
            flags |= FLAG_FIRST
        if i == total - 1:
            flags |= FLAG_LAST
        crc = zlib.crc32(data) & 0xFFFFFFFF
        header = _make_header(
            seq=i, total=total, crc=crc, flags=flags
        )
        chunks.append(header + data)

    return chunks
\end{lstlisting}

\newpage
\subsection*{В.3 Модуль реассемблирования (\texttt{core/reassembly.py})}

Модуль реализует сборку фрагментированных данных с поддержкой
получения чанков в произвольном порядке, дедупликацией избыточных
копий, проверкой CRC32 каждого чанка и верификацией SHA-256
собранного файла.

\begin{lstlisting}[style=appendixcode]
"""Chunk reassembly with out-of-order support and
integrity verification.

Collects fragmented chunks, reorders by sequence
number, verifies CRC32 per chunk and SHA-256 of
the full payload.
"""

import hashlib
import zlib

from netstego.core.fragmentation import (
    FLAG_FIRST, FLAG_LAST, parse_header,
)


class ReassemblyError(Exception):
    """Raised when reassembly fails."""


class Reassembler:
    """Collects chunks and reassembles the original
    payload. Supports out-of-order delivery."""

    def __init__(self) -> None:
        self._chunks: dict[int, tuple[int, bytes]] = {}
        self._total: int | None = None
        self._duplicates: int = 0

    @property
    def received_count(self) -> int:
        """Number of unique chunks received."""
        return len(self._chunks)

    @property
    def total_expected(self) -> int | None:
        """Total chunks expected, or None if unknown."""
        return self._total

    @property
    def missing_seqs(self) -> list[int]:
        """List of missing sequence numbers."""
        if self._total is None:
            return []
        return [
            i for i in range(self._total)
            if i not in self._chunks
        ]

    @property
    def duplicates_received(self) -> int:
        """Number of duplicate chunks discarded."""
        return self._duplicates

    def add_chunk(self, raw_chunk: bytes) -> int:
        """Parse and store a chunk.

        Duplicate chunks are silently ignored, enabling
        redundancy-based packet loss tolerance.

        Raises:
            ValueError: If chunk is malformed.
            ReassemblyError: If CRC32 check fails.
        """
        seq, total, crc, flags, data = parse_header(
            raw_chunk
        )

        actual_crc = zlib.crc32(data) & 0xFFFFFFFF
        if actual_crc != crc:
            msg = (
                f"CRC32 mismatch for chunk {seq}: "
                f"expected {crc:#010x}, "
                f"got {actual_crc:#010x}"
            )
            raise ReassemblyError(msg)

        if self._total is None:
            self._total = total
        elif self._total != total:
            msg = (
                f"Total mismatch: expected "
                f"{self._total}, got {total}"
            )
            raise ReassemblyError(msg)

        if seq in self._chunks:
            self._duplicates += 1
            return seq

        self._chunks[seq] = (flags, data)
        return seq

    def is_complete(self) -> bool:
        """Check if all chunks have been received."""
        if self._total is None:
            return False
        return len(self._chunks) == self._total

    def try_reassemble(self) -> bytes:
        """Attempt to reassemble the full payload.

        Raises:
            ReassemblyError: If chunks are missing,
                flags inconsistent, or SHA-256 fails.
        """
        if not self.is_complete():
            missing = self.missing_seqs
            msg = (
                f"Missing {len(missing)} chunks: "
                f"{missing[:10]}"
            )
            raise ReassemblyError(msg)

        ordered = sorted(self._chunks.items())

        first_flags = ordered[0][1][0]
        if not (first_flags & FLAG_FIRST):
            msg = "First chunk missing FLAG_FIRST"
            raise ReassemblyError(msg)

        last_flags = ordered[-1][1][0]
        if not (last_flags & FLAG_LAST):
            msg = "Last chunk missing FLAG_LAST"
            raise ReassemblyError(msg)

        first_data = ordered[0][1][1]
        if len(first_data) < 32:
            msg = (
                "First chunk too short to contain "
                "SHA-256 hash"
            )
            raise ReassemblyError(msg)

        expected_hash = first_data[:32]
        payload_parts = [first_data[32:]]

        for _seq, (_, data) in ordered[1:]:
            payload_parts.append(data)

        payload = b"".join(payload_parts)

        actual_hash = hashlib.sha256(payload).digest()
        if actual_hash != expected_hash:
            msg = (
                "SHA-256 verification failed: "
                "payload corrupted"
            )
            raise ReassemblyError(msg)

        return payload

    def reset(self) -> None:
        """Clear all received chunks."""
        self._chunks.clear()
        self._total = None
        self._duplicates = 0
\end{lstlisting}

\newpage
\section*{ПРИЛОЖЕНИЕ Г. Реализация стеганографических каналов}
\addcontentsline{toc}{section}{ПРИЛОЖЕНИЕ Г. Реализация стеганографических каналов}

\subsection*{Г.1 Базовый класс каналов (\texttt{channels/base.py})}

\begin{lstlisting}[style=appendixcode]
"""Abstract base class for steganographic channels.

Every channel implementation must inherit from
StegoChannel and implement all abstract methods.
"""

from abc import ABC, abstractmethod
from scapy.packet import Packet


class StegoChannel(ABC):
    """Base class for all steganographic channels."""

    @abstractmethod
    def encode(
        self, data: bytes, dest_ip: str, **kwargs,
    ) -> list[Packet]:
        """Encode data into a list of Scapy packets."""

    @abstractmethod
    def decode(
        self, packets: list[Packet],
    ) -> bytes:
        """Decode hidden data from a list of packets."""

    @property
    @abstractmethod
    def bits_per_packet(self) -> int:
        """Number of data bits per single packet."""

    @property
    @abstractmethod
    def protocol_filter(self) -> str:
        """BPF filter string for the receiver sniffer.
        May contain {ip} placeholder for source IP."""

    @property
    def name(self) -> str:
        """Human-readable channel name."""
        return self.__class__.__name__
\end{lstlisting}

\subsection*{Г.2 Канал ICMP Payload (\texttt{channels/icmp\_payload.py})}

Канал использует поле данных ICMP Echo Request. Максимальная ёмкость ---
1472~байт на пакет (MTU~1500 $-$ 20~IP $-$ 8~ICMP).

\begin{lstlisting}[style=appendixcode]
"""ICMP Echo payload steganographic channel.

Hides data in the payload (data) field of ICMP Echo
Request packets. Max payload: 1472 bytes
(MTU 1500 - 20 IP header - 8 ICMP header).

Reference: RFC 792.
"""

import struct
from scapy.layers.inet import ICMP, IP
from scapy.packet import Packet, Raw
from netstego.channels.base import StegoChannel

MAX_ICMP_PAYLOAD = 1472


class IcmpPayloadChannel(StegoChannel):
    """Steganographic channel using ICMP Echo payload.

    Each packet carries up to 1472 bytes of hidden data.
    A 2-byte length prefix is prepended so the decoder
    knows how many bytes are valid in each packet.
    """

    @property
    def bits_per_packet(self) -> int:
        return MAX_ICMP_PAYLOAD * 8

    @property
    def protocol_filter(self) -> str:
        return "icmp and src host {ip}"

    def encode(
        self, data: bytes, dest_ip: str, **kwargs,
    ) -> list[Packet]:
        """Encode data into ICMP Echo Request packets."""
        max_data = MAX_ICMP_PAYLOAD - 2
        packets: list[Packet] = []
        offset = 0
        seq = 0

        while offset < len(data):
            chunk = data[offset : offset + max_data]
            payload = (
                struct.pack(">H", len(chunk)) + chunk
            )
            pkt = (
                IP(dst=dest_ip)
                / ICMP(type=8, code=0, seq=seq)
                / Raw(load=payload)
            )
            packets.append(pkt)
            offset += max_data
            seq += 1

        return packets

    def decode(
        self, packets: list[Packet],
    ) -> bytes:
        """Decode hidden data from ICMP packets."""
        sorted_pkts = sorted(
            packets, key=lambda p: p[ICMP].seq
        )
        parts: list[bytes] = []

        for pkt in sorted_pkts:
            raw = bytes(pkt[Raw].load)
            length = struct.unpack(">H", raw[:2])[0]
            parts.append(raw[2 : 2 + length])

        return b"".join(parts)
\end{lstlisting}

\newpage
\subsection*{Г.3 Канал IP Identification (\texttt{channels/ip\_id.py})}

Канал встраивает данные в 16-битное поле идентификации IPv4-заголовка.
Ёмкость --- 2~байта на пакет.

\begin{lstlisting}[style=appendixcode]
"""IP Identification field steganographic channel.

Hides data in the 16-bit IP ID field of IPv4 packets.
Each packet carries 2 bytes (16 bits) of hidden data.

Reference: RFC 791, Section 3.1.
"""

import struct
from scapy.layers.inet import IP, ICMP
from scapy.packet import Packet, Raw
from netstego.channels.base import StegoChannel

BITS_PER_PACKET = 16


class IpIdChannel(StegoChannel):
    """Steganographic channel using IPv4 ID field.

    Each packet embeds 2 bytes of data in IP ID.
    ICMP Echo Request is used as carrier protocol.
    """

    @property
    def bits_per_packet(self) -> int:
        return BITS_PER_PACKET

    @property
    def protocol_filter(self) -> str:
        return "icmp and src host {ip}"

    def encode(
        self, data: bytes, dest_ip: str, **kwargs,
    ) -> list[Packet]:
        """Encode data into IP packets using ID field."""
        packets: list[Packet] = []
        offset = 0
        seq = 0

        while offset < len(data):
            chunk = data[offset : offset + 2]
            if len(chunk) == 1:
                chunk = chunk + b"\x00"
            ip_id = struct.unpack(">H", chunk)[0]
            pkt = (
                IP(dst=dest_ip, id=ip_id)
                / ICMP(type=8, code=0, seq=seq)
                / Raw(load=b"\x00" * 32)
            )
            packets.append(pkt)
            offset += 2
            seq += 1

        # Termination packet with total data length
        term_payload = struct.pack(">I", len(data))
        pkt = (
            IP(dst=dest_ip, id=0xFFFF)
            / ICMP(type=8, code=0, seq=seq)
            / Raw(load=term_payload)
        )
        packets.append(pkt)
        return packets

    def decode(
        self, packets: list[Packet],
    ) -> bytes:
        """Decode hidden data from IP ID fields."""
        icmp_pkts = [
            p for p in packets if p.haslayer(ICMP)
        ]
        sorted_pkts = sorted(
            icmp_pkts, key=lambda p: p[ICMP].seq
        )

        if not sorted_pkts:
            return b""

        term_pkt = sorted_pkts[-1]
        term_raw = (
            bytes(term_pkt[Raw].load)
            if term_pkt.haslayer(Raw)
            else b""
        )
        if (
            len(term_raw) >= 4
            and term_pkt[IP].id == 0xFFFF
        ):
            total_len = struct.unpack(
                ">I", term_raw[:4]
            )[0]
            data_pkts = sorted_pkts[:-1]
        else:
            total_len = len(sorted_pkts) * 2
            data_pkts = sorted_pkts

        parts: list[bytes] = []
        for pkt in data_pkts:
            ip_id = pkt[IP].id
            parts.append(struct.pack(">H", ip_id))

        result = b"".join(parts)
        return result[:total_len]
\end{lstlisting}

\newpage
\subsection*{Г.4 Канал TCP ISN (\texttt{channels/tcp\_isn.py})}

Канал встраивает данные в 32-битное поле начального порядкового номера
TCP-сегмента (SYN-пакеты). Ёмкость --- 4~байта на пакет.

\begin{lstlisting}[style=appendixcode]
"""TCP Initial Sequence Number steganographic channel.

Hides data in the 32-bit TCP Sequence Number field
of SYN packets. Each SYN carries 4 bytes of data.

Reference: RFC 793, Section 3.1.
           RFC 6528 -- Sequence Number Attacks.
"""

import struct
from scapy.layers.inet import IP, TCP
from scapy.packet import Packet
from netstego.channels.base import StegoChannel

BITS_PER_PACKET = 32


class TcpIsnChannel(StegoChannel):
    """Steganographic channel using TCP ISN.

    Each TCP SYN embeds 4 bytes of data in seq number.
    Uses incrementing destination ports for ordering.
    """

    @property
    def bits_per_packet(self) -> int:
        return BITS_PER_PACKET

    @property
    def protocol_filter(self) -> str:
        return (
            "tcp[tcpflags] & tcp-syn != 0 "
            "and src host {ip}"
        )

    def encode(
        self, data: bytes, dest_ip: str, **kwargs,
    ) -> list[Packet]:
        """Encode data into TCP SYN packets."""
        base_port = kwargs.get("base_port", 10000)
        packets: list[Packet] = []
        offset = 0
        port_offset = 0

        while offset < len(data):
            chunk = data[offset : offset + 4]
            padded = chunk.ljust(4, b"\x00")
            isn = struct.unpack(">I", padded)[0]
            dport = base_port + port_offset
            pkt = IP(dst=dest_ip) / TCP(
                sport=40000 + port_offset,
                dport=dport,
                seq=isn,
                flags="S",
            )
            packets.append(pkt)
            offset += 4
            port_offset += 1

        # Termination SYN with data length in seq
        term_pkt = IP(dst=dest_ip) / TCP(
            sport=65535,
            dport=base_port + port_offset,
            seq=len(data),
            flags="S",
        )
        packets.append(term_pkt)
        return packets

    def decode(
        self, packets: list[Packet],
    ) -> bytes:
        """Decode hidden data from TCP SYN ISNs."""
        tcp_pkts = [
            p for p in packets
            if p.haslayer(TCP) and p[TCP].flags & 0x02
        ]
        sorted_pkts = sorted(
            tcp_pkts, key=lambda p: p[TCP].dport
        )

        if not sorted_pkts:
            return b""

        term_pkt = sorted_pkts[-1]
        if term_pkt[TCP].sport == 65535:
            total_len = term_pkt[TCP].seq
            data_pkts = sorted_pkts[:-1]
        else:
            total_len = len(sorted_pkts) * 4
            data_pkts = sorted_pkts

        parts: list[bytes] = []
        for pkt in data_pkts:
            isn = pkt[TCP].seq
            parts.append(struct.pack(">I", isn))

        result = b"".join(parts)
        return result[:total_len]
\end{lstlisting}

\newpage
\subsection*{Г.5 Канал DNS QNAME (\texttt{channels/dns\_query.py})}

Канал кодирует данные в поддоменных метках DNS-запросов с использованием
Base32. Ёмкость --- до 30~байт на пакет. Суффикс домена является
конфигурируемым параметром.

\begin{lstlisting}[style=appendixcode]
"""DNS QNAME subdomain steganographic channel.

Hides data in DNS query names using base32-encoded
subdomains. Each DNS query carries up to ~30 raw
bytes per packet.

Reference: RFC 1035, Section 4.
"""

import base64
import struct
from scapy.layers.dns import DNS, DNSQR
from scapy.layers.inet import IP, UDP
from scapy.packet import Packet
from netstego.channels.base import StegoChannel

MAX_RAW_BYTES = 30
DEFAULT_DOMAIN_SUFFIX = "cdn-cache.net"


class DnsQueryChannel(StegoChannel):
    """Steganographic channel using DNS QNAME.

    Data is base32-encoded into subdomain labels.
    The domain suffix is configurable to avoid
    trivial signature detection.
    """

    def __init__(
        self, domain_suffix: str | None = None,
    ) -> None:
        self._suffix = (
            domain_suffix or DEFAULT_DOMAIN_SUFFIX
        )

    @property
    def domain_suffix(self) -> str:
        return self._suffix

    @property
    def bits_per_packet(self) -> int:
        return MAX_RAW_BYTES * 8

    @property
    def protocol_filter(self) -> str:
        return "udp port 53 and src host {ip}"

    def encode(
        self, data: bytes, dest_ip: str, **kwargs,
    ) -> list[Packet]:
        """Encode data into DNS query packets."""
        packets: list[Packet] = []
        offset = 0
        tx_id = 0

        while offset < len(data):
            chunk = data[offset : offset + MAX_RAW_BYTES]
            encoded = (
                base64.b32encode(chunk)
                .rstrip(b"=")
                .decode("ascii")
                .lower()
            )
            labels = []
            for i in range(0, len(encoded), 63):
                labels.append(encoded[i : i + 63])
            qname = (
                ".".join(labels)
                + "." + self._suffix + "."
            )

            pkt = (
                IP(dst=dest_ip)
                / UDP(sport=30000 + tx_id, dport=53)
                / DNS(
                    id=tx_id,
                    qd=DNSQR(qname=qname, qtype="A"),
                )
            )
            packets.append(pkt)
            offset += MAX_RAW_BYTES
            tx_id += 1

        # Termination query with data length
        term_payload = struct.pack(">I", len(data))
        term_encoded = (
            base64.b32encode(term_payload)
            .rstrip(b"=")
            .decode("ascii")
            .lower()
        )
        term_qname = (
            f"{term_encoded}.term.{self._suffix}."
        )
        pkt = (
            IP(dst=dest_ip)
            / UDP(sport=30000 + tx_id, dport=53)
            / DNS(
                id=tx_id,
                qd=DNSQR(qname=term_qname, qtype="A"),
            )
        )
        packets.append(pkt)
        return packets

    def decode(
        self, packets: list[Packet],
    ) -> bytes:
        """Decode hidden data from DNS query names."""
        dns_pkts = [
            p for p in packets
            if p.haslayer(DNS) and p.haslayer(DNSQR)
        ]
        sorted_pkts = sorted(
            dns_pkts, key=lambda p: p[DNS].id
        )

        if not sorted_pkts:
            return b""

        total_len = None
        data_pkts = []

        for pkt in sorted_pkts:
            qname = pkt[DNSQR].qname
            if isinstance(qname, bytes):
                qname = qname.decode(
                    "ascii", errors="ignore"
                )
            qname = qname.rstrip(".")

            parts = qname.split(".")

            if (
                len(parts) >= 3
                and parts[-2] == "term"
            ):
                encoded = parts[0].upper()
                pad = (8 - len(encoded) % 8) % 8
                encoded += "=" * pad
                try:
                    length_bytes = base64.b32decode(
                        encoded
                    )
                    if len(length_bytes) >= 4:
                        total_len = struct.unpack(
                            ">I", length_bytes[:4]
                        )[0]
                except Exception:
                    pass
            else:
                data_pkts.append(pkt)

        parts_data: list[bytes] = []
        for pkt in data_pkts:
            qname = pkt[DNSQR].qname
            if isinstance(qname, bytes):
                qname = qname.decode(
                    "ascii", errors="ignore"
                )
            qname = qname.rstrip(".")

            labels = qname.split(".")
            data_labels = labels[:-2]
            encoded = "".join(data_labels).upper()
            pad = (8 - len(encoded) % 8) % 8
            encoded += "=" * pad
            try:
                decoded = base64.b32decode(encoded)
                parts_data.append(decoded)
            except Exception:
                continue

        result = b"".join(parts_data)
        if total_len is not None:
            result = result[:total_len]
        return result
\end{lstlisting}

\newpage
\subsection*{Г.6 Канал TCP Timestamp (\texttt{channels/tcp\_timestamp.py})}

Канал встраивает данные в поле TSval опции TCP Timestamp.
Ёмкость --- 4~байта (32~бита) на пакет.

\begin{lstlisting}[style=appendixcode]
"""TCP Timestamp steganographic channel.

Hides data in the TCP Timestamp Value (TSval) option.
Each packet carries 4 bytes (32 bits) of hidden data.

Reference: RFC 7323, Section 3.
"""

import struct
from scapy.layers.inet import IP, TCP
from scapy.packet import Packet
from netstego.channels.base import StegoChannel

BITS_PER_PACKET = 32


class TcpTimestampChannel(StegoChannel):
    """Steganographic channel using TCP TSval.

    Each packet embeds 4 bytes of data in the TCP
    Timestamp Value option. Uses ACK packets with
    incrementing sequence numbers for ordering.
    """

    @property
    def bits_per_packet(self) -> int:
        return BITS_PER_PACKET

    @property
    def protocol_filter(self) -> str:
        return "tcp and src host {ip}"

    def encode(
        self, data: bytes, dest_ip: str, **kwargs,
    ) -> list[Packet]:
        """Encode data into TCP packets using TSval."""
        base_port = kwargs.get("base_port", 20000)
        packets: list[Packet] = []
        offset = 0
        seq_num = 0

        while offset < len(data):
            chunk = data[offset : offset + 4]
            padded = chunk.ljust(4, b"\x00")
            tsval = struct.unpack(">I", padded)[0]
            ts_option = ("Timestamp", (tsval, 0))
            pkt = IP(dst=dest_ip) / TCP(
                sport=50000 + seq_num,
                dport=base_port,
                seq=seq_num,
                flags="A",
                options=[ts_option],
            )
            packets.append(pkt)
            offset += 4
            seq_num += 1

        # Termination: TSecr=1 marks termination
        ts_option = ("Timestamp", (len(data), 1))
        term_pkt = IP(dst=dest_ip) / TCP(
            sport=50000 + seq_num,
            dport=base_port,
            seq=seq_num,
            flags="A",
            options=[ts_option],
        )
        packets.append(term_pkt)
        return packets

    def decode(
        self, packets: list[Packet],
    ) -> bytes:
        """Decode hidden data from TCP Timestamp."""
        tcp_pkts = [
            p for p in packets if p.haslayer(TCP)
        ]
        sorted_pkts = sorted(
            tcp_pkts, key=lambda p: p[TCP].seq
        )

        if not sorted_pkts:
            return b""

        total_len = None
        data_pkts = []

        for pkt in sorted_pkts:
            options = pkt[TCP].options
            tsval, tsecr = None, None
            for opt_name, opt_val in options:
                if opt_name == "Timestamp":
                    tsval, tsecr = opt_val
                    break

            if tsval is None:
                continue

            if tsecr == 1:
                total_len = tsval
            else:
                data_pkts.append((tsval,))

        parts: list[bytes] = []
        for (tsval,) in data_pkts:
            parts.append(struct.pack(">I", tsval))

        result = b"".join(parts)
        if total_len is not None:
            result = result[:total_len]
        return result
\end{lstlisting}

\newpage
\section*{ПРИЛОЖЕНИЕ Д. Сетевой конвейер отправки и приёма данных}
\addcontentsline{toc}{section}{ПРИЛОЖЕНИЕ Д. Сетевой конвейер отправки и приёма данных}

\subsection*{Д.1 Модуль отправки (\texttt{network/sender.py})}

Модуль реализует полный конвейер отправки: чтение файла, шифрование
AES-256-GCM, фрагментацию с NST-заголовками, кодирование через
выбранный стеганографический канал, отправку с управлением
скоростью и джиттером, сбор статистики и сохранение в PCAP.

\begin{lstlisting}[style=appendixcode]
"""Sender module: encrypts, fragments, encodes,
and transmits steganographic packets.

Handles the full pipeline:
    1. Read file -> encrypt with AES-256-GCM
    2. Fragment ciphertext into chunks with headers
    3. Frame chunks with length prefixes and encode
       via StegoChannel
    4. Send packets with configurable rate and jitter
    5. Optionally save sent packets to PCAP
"""

import platform
import random
import struct
import sys
import time
from pathlib import Path

from rich.console import Console
from rich.progress import Progress
from scapy.sendrecv import send as scapy_send
from scapy.utils import wrpcap

from netstego.channels.base import StegoChannel
from netstego.config import TimingConfig
from netstego.core.crypto import encrypt, load_key
from netstego.core.fragmentation import fragment
from netstego.stats.collector import StatsCollector

console = Console()


def check_privileges() -> None:
    """Verify raw socket privileges.

    On Linux/macOS checks for root or CAP_NET_RAW.
    On Windows requires Administrator.
    """
    if platform.system() == "Windows":
        import ctypes
        if not ctypes.windll.shell32.IsUserAnAdmin():
            console.print(
                "[red]Error:[/red] "
                "Administrator privileges required."
            )
            sys.exit(1)
    else:
        import os
        if os.geteuid() != 0:
            try:
                import socket
                s = socket.socket(
                    socket.AF_PACKET,
                    socket.SOCK_RAW,
                )
                s.close()
            except (PermissionError, OSError):
                console.print(
                    "[red]Error:[/red] "
                    "Root or CAP_NET_RAW required."
                )
                sys.exit(1)


def get_channel(
    channel_name: str, **kwargs,
) -> StegoChannel:
    """Instantiate a StegoChannel by config name.

    Args:
        channel_name: One of 'icmp', 'ip-id',
            'tcp-isn', 'dns', 'tcp-ts'.
        **kwargs: Channel-specific options.

    Raises:
        ValueError: If channel name is unknown.
    """
    if channel_name == "icmp":
        from netstego.channels.icmp_payload import (
            IcmpPayloadChannel,
        )
        return IcmpPayloadChannel()
    if channel_name == "ip-id":
        from netstego.channels.ip_id import (
            IpIdChannel,
        )
        return IpIdChannel()
    if channel_name == "tcp-isn":
        from netstego.channels.tcp_isn import (
            TcpIsnChannel,
        )
        return TcpIsnChannel()
    if channel_name == "dns":
        from netstego.channels.dns_query import (
            DnsQueryChannel,
        )
        return DnsQueryChannel(
            domain_suffix=kwargs.get("domain_suffix")
        )
    if channel_name == "tcp-ts":
        from netstego.channels.tcp_timestamp import (
            TcpTimestampChannel,
        )
        return TcpTimestampChannel()

    msg = f"Unknown channel: {channel_name!r}"
    raise ValueError(msg)


def frame_chunks(chunks: list[bytes]) -> bytes:
    """Frame chunks with 4-byte length prefixes.

    Format: [len1(4)][chunk1][len2(4)][chunk2]...
    """
    parts = []
    for chunk in chunks:
        parts.append(struct.pack(">I", len(chunk)))
        parts.append(chunk)
    return b"".join(parts)


def send_data(
    data: bytes,
    dest_ip: str,
    key_path: Path,
    channel_name: str,
    timing: TimingConfig | None = None,
    redundancy: int = 1,
    pcap_out: Path | None = None,
    db_path: str | None = None,
) -> int:
    """Encrypt, fragment, and send data through
    a steganographic channel.

    Args:
        data: Raw plaintext bytes to send.
        dest_ip: Destination IP address.
        key_path: Path to the AES-256 key file.
        channel_name: Steganographic channel name.
        timing: Timing config (rate, jitter).
        redundancy: Chunk duplication factor.
        pcap_out: Path to save packets as PCAP.
        db_path: Path to stats database.

    Returns:
        Total number of packets sent.
    """
    check_privileges()

    if timing is None:
        timing = TimingConfig()

    key = load_key(key_path)
    console.print(
        f"[dim]Key loaded from {key_path}[/dim]"
    )
    console.print(
        f"[dim]Data size: {len(data)} bytes[/dim]"
    )

    ciphertext = encrypt(data, key)
    console.print(
        f"[dim]Ciphertext size: "
        f"{len(ciphertext)} bytes[/dim]"
    )

    channel = get_channel(channel_name)

    bytes_per_packet = channel.bits_per_packet // 8
    overhead = 17 + 4
    if bytes_per_packet > overhead + 33:
        chunk_data_size = bytes_per_packet - overhead
    else:
        chunk_data_size = 128

    chunks = fragment(ciphertext, chunk_data_size)
    console.print(
        f"[dim]Fragmented into {len(chunks)} chunks "
        f"({chunk_data_size} data bytes/chunk)[/dim]"
    )

    if redundancy > 1:
        redundant_chunks = []
        for chunk in chunks:
            for _ in range(redundancy):
                redundant_chunks.append(chunk)
        random.shuffle(redundant_chunks)
        console.print(
            f"[dim]Redundancy x{redundancy}: "
            f"{len(redundant_chunks)} chunks "
            f"({len(chunks)} unique)[/dim]"
        )
        chunks = redundant_chunks

    framed = frame_chunks(chunks)
    all_packets = channel.encode(framed, dest_ip)

    total_packets = len(all_packets)
    console.print(
        f"[bold green]Sending {total_packets} "
        f"packets[/bold green] via "
        f"[cyan]{channel.name}[/cyan] to "
        f"[cyan]{dest_ip}[/cyan]"
    )

    collector = StatsCollector(db_path=db_path)
    session_id = collector.start_session(
        direction="send",
        channel=channel_name,
        dest_ip=dest_ip,
        file_size=len(data),
        total_packets=total_packets,
        total_chunks=len(chunks),
    )

    base_delay = (
        1.0 / timing.rate_pps
        if timing.rate_pps > 0
        else 0.0
    )

    try:
        with Progress(console=console) as progress:
            task = progress.add_task(
                "[green]Transmitting...",
                total=total_packets,
            )
            for idx, pkt in enumerate(all_packets):
                scapy_send(pkt, verbose=False)
                collector.log_packet(
                    session_id, seq=idx,
                    field_name=channel.name,
                    size=len(pkt),
                )
                progress.advance(task)

                if base_delay > 0:
                    jitter = random.uniform(
                        -timing.jitter_factor
                        * base_delay,
                        timing.jitter_factor
                        * base_delay,
                    )
                    delay = max(
                        0.0, base_delay + jitter
                    )
                    time.sleep(delay)

        collector.finish_session(
            session_id, status="completed",
            total_packets=total_packets,
        )
    except Exception:
        collector.finish_session(
            session_id, status="failed",
            total_packets=total_packets,
        )
        raise

    if pcap_out is not None:
        wrpcap(str(pcap_out), all_packets)
        console.print(
            f"[dim]Packets saved to {pcap_out}[/dim]"
        )

    console.print(
        f"[bold green]Done![/bold green] Sent "
        f"{total_packets} packets "
        f"({len(ciphertext)} bytes encrypted). "
        f"Session: {session_id}"
    )
    return total_packets


def send_file(
    file_path: Path,
    dest_ip: str,
    key_path: Path,
    channel_name: str,
    timing: TimingConfig | None = None,
    redundancy: int = 1,
    pcap_out: Path | None = None,
    db_path: str | None = None,
) -> int:
    """Read a file and send via send_data."""
    plaintext = file_path.read_bytes()
    return send_data(
        data=plaintext,
        dest_ip=dest_ip,
        key_path=key_path,
        channel_name=channel_name,
        timing=timing,
        redundancy=redundancy,
        pcap_out=pcap_out,
        db_path=db_path,
    )
\end{lstlisting}

\newpage
\subsection*{Д.2 Модуль приёма (\texttt{network/receiver.py})}

Модуль реализует полный конвейер приёма: перехват пакетов (BPF-фильтр
или чтение из PCAP), декодирование через стеганографический канал,
извлечение и реассемблирование чанков, расшифрование AES-256-GCM
и запись результата в файл.

\begin{lstlisting}[style=appendixcode]
"""Receiver module: sniffs packets, decodes,
reassembles, and decrypts.

Handles the full receive pipeline:
    1. Sniff packets via AsyncSniffer with BPF filter
       (or read from PCAP)
    2. Decode packets via the chosen StegoChannel
    3. Unframe length-prefixed chunks
    4. Reassemble chunks with out-of-order support
    5. Decrypt with AES-256-GCM and write output file
"""

import struct
import time as _time
from pathlib import Path

from rich.console import Console
from scapy.sendrecv import AsyncSniffer
from scapy.utils import rdpcap, wrpcap

from netstego.core.crypto import decrypt, load_key
from netstego.core.reassembly import Reassembler
from netstego.network.sender import (
    check_privileges, get_channel,
)
from netstego.stats.collector import StatsCollector

console = Console()


def unframe_chunks(data: bytes) -> list[bytes]:
    """Extract length-prefixed chunks from a framed
    byte stream.

    Format: [len1(4)][chunk1][len2(4)][chunk2]...
    """
    chunks = []
    offset = 0
    while offset + 4 <= len(data):
        length = struct.unpack(
            ">I", data[offset : offset + 4]
        )[0]
        offset += 4
        if offset + length > len(data):
            break
        chunks.append(data[offset : offset + length])
        offset += length
    return chunks


def receive_file(
    sender_ip: str,
    key_path: Path,
    output_path: Path,
    channel_name: str,
    timeout: int = 60,
    pcap_out: Path | None = None,
    pcap_in: Path | None = None,
    iface: str | None = None,
    db_path: str | None = None,
) -> bool:
    """Receive and reconstruct a file from a
    steganographic channel.

    Args:
        sender_ip: IP address of the sender.
        key_path: Path to the AES-256 key file.
        output_path: Path to write decrypted output.
        channel_name: Steganographic channel name.
        timeout: Sniffing timeout in seconds.
        pcap_out: Path to save captured packets.
        pcap_in: PCAP file to read instead of sniffing.
        iface: Network interface name for sniffing.
        db_path: Path to stats database.

    Returns:
        True if file was successfully received.
    """
    if pcap_in is None:
        check_privileges()

    key = load_key(key_path)
    console.print(
        f"[dim]Key loaded from {key_path}[/dim]"
    )

    channel = get_channel(channel_name)

    collector = StatsCollector(db_path=db_path)
    session_id = collector.start_session(
        direction="receive",
        channel=channel_name,
        bind_ip=sender_ip,
    )

    if pcap_in is not None:
        console.print(
            f"[bold blue]Reading from PCAP:[/bold blue]"
            f" {pcap_in}"
        )
        all_packets = list(rdpcap(str(pcap_in)))
    else:
        bpf_filter = channel.protocol_filter.format(
            ip=sender_ip
        )
        console.print(
            f"[bold blue]Listening[/bold blue] for "
            f"packets from [cyan]{sender_ip}[/cyan] "
            f"via [cyan]{channel.name}[/cyan] "
            f"(timeout={timeout}s)"
        )
        console.print(
            f"[dim]BPF filter: {bpf_filter}[/dim]"
        )
        if iface:
            console.print(
                f"[dim]Interface: {iface}[/dim]"
            )

        sniffer_kwargs = {
            "filter": bpf_filter, "store": True,
        }
        if iface:
            sniffer_kwargs["iface"] = iface

        sniffer = AsyncSniffer(**sniffer_kwargs)
        sniffer.start()

        try:
            _time.sleep(timeout)
        except KeyboardInterrupt:
            console.print(
                "\n[yellow]Interrupted[/yellow]"
            )
        finally:
            sniffer.stop()

        all_packets = (
            list(sniffer.results)
            if sniffer.results
            else []
        )

    console.print(
        f"[dim]Captured {len(all_packets)} "
        f"packets[/dim]"
    )

    for idx, pkt in enumerate(all_packets):
        collector.log_packet(
            session_id, seq=idx,
            field_name=channel.name,
            size=len(pkt),
        )

    if not all_packets:
        console.print(
            "[red]Error:[/red] No packets received."
        )
        collector.finish_session(
            session_id, status="failed"
        )
        return False

    if pcap_out is not None:
        wrpcap(str(pcap_out), all_packets)
        console.print(
            f"[dim]Packets saved to {pcap_out}[/dim]"
        )

    try:
        raw_data = channel.decode(all_packets)
    except Exception as e:
        console.print(
            f"[red]Error decoding:[/red] {e}"
        )
        collector.finish_session(
            session_id, status="failed",
            total_packets=len(all_packets),
        )
        return False

    console.print(
        f"[dim]Decoded {len(raw_data)} bytes[/dim]"
    )

    chunks = unframe_chunks(raw_data)
    console.print(
        f"[dim]Extracted {len(chunks)} chunks[/dim]"
    )

    if not chunks:
        console.print(
            "[red]Error:[/red] No valid chunks found."
        )
        collector.finish_session(
            session_id, status="failed",
            total_packets=len(all_packets),
        )
        return False

    reassembler = Reassembler()
    for chunk in chunks:
        try:
            reassembler.add_chunk(chunk)
        except Exception as e:
            console.print(
                f"[yellow]Warning:[/yellow] "
                f"Skipping bad chunk: {e}"
            )

    if reassembler.duplicates_received > 0:
        console.print(
            f"[dim]Deduplicated "
            f"{reassembler.duplicates_received} "
            f"redundant chunks[/dim]"
        )

    if not reassembler.is_complete():
        missing = reassembler.missing_seqs
        total = reassembler.total_expected or 0
        received = reassembler.received_count
        loss_pct = (
            (len(missing) / total * 100)
            if total > 0
            else 0
        )
        console.print(
            f"[red]Error:[/red] Reassembly incomplete."
            f" Received {received}/{total} chunks "
            f"({loss_pct:.1f}% loss). "
            f"Missing: {missing[:10]}"
        )
        collector.finish_session(
            session_id, status="failed",
            total_packets=len(all_packets),
        )
        return False

    try:
        ciphertext = reassembler.try_reassemble()
    except Exception as e:
        console.print(
            f"[red]Error during reassembly:[/red] {e}"
        )
        collector.finish_session(
            session_id, status="failed",
            total_packets=len(all_packets),
        )
        return False

    console.print(
        f"[dim]Reassembled {len(ciphertext)} bytes "
        f"of ciphertext[/dim]"
    )

    try:
        plaintext = decrypt(ciphertext, key)
    except Exception as e:
        console.print(
            f"[red]Error during decryption:[/red] {e}"
        )
        collector.finish_session(
            session_id, status="failed",
            total_packets=len(all_packets),
        )
        return False

    output_path.parent.mkdir(
        parents=True, exist_ok=True
    )
    output_path.write_bytes(plaintext)

    collector.finish_session(
        session_id, status="completed",
        total_packets=len(all_packets),
    )
    console.print(
        f"[bold green]Done![/bold green] Received "
        f"{len(plaintext)} bytes -> {output_path} "
        f"(session: {session_id})"
    )
    return True
\end{lstlisting}
